\section{Introduction}
\label{sec:introduction}

Semantic segmentation on heterogeneous edge accelerators is usually evaluated through a proxy such as FLOPs or parameter count, even though those proxies need not predict measured latency or energy on different accelerator architectures. A graph compiled for one backend is also not guaranteed to compile, unmodified, for another backend. Finally, clean-image mIoU does not describe reliability under night, rain, fog, or snow. Training one independently optimized model per device can further increase training and maintenance cost.

We study the underexplored intersection of compiler-portable elastic segmentation, hardware-measured resource optimization, and calibrated risk-aware inference under adverse visual conditions. The working hypothesis is that a shared, compiler-safe supernet can provide several static deployable choices while preserving accuracy close to an independently trained model of comparable capacity.

The intended contributions are:
\begin{enumerate}
  \item a restricted-operator elastic segmentation supernet with static subnets for heterogeneous compilers;
  \item a hardware-in-the-loop Pareto search that joins measured backend latency with segmentation quality;
  \item a calibrated visual-risk router that can select among precompiled engines; and
  \item a cross-platform reliability and measurement protocol spanning clean and adverse conditions.
\end{enumerate}

The first two contributions have partial evidence in the current repository. The compiler-safe supernet, training loop, evaluation loop, benchmark harness, and first Pareto lookup-table search are implemented. The risk router, quantization-aware training, energy axis, systematic literature review, and the full ablation matrix remain open.

\TODO{Replace this paragraph with a final novelty statement only after the systematic literature review is complete.}
